\documentclass[11pt, a4paper]{article}

\usepackage[a4paper, top=2.2cm, bottom=2.2cm, left=2.4cm, right=2.4cm]{geometry}
\usepackage{fontenc}
\usepackage{inputenc}
\usepackage{microtype}
\usepackage{parskip}
\usepackage{booktabs}
\usepackage{tabularx}
\usepackage{array}
\usepackage{enumitem}
\usepackage{titlesec}
\usepackage{fancyhdr}
\usepackage{hyperref}
\usepackage{mdframed}
\usepackage{setspace}

\hypersetup{
  colorlinks=false,
  hidelinks
}

% Typography — serif font for direct Overleaf compilation
\renewcommand{\familydefault}{\rmdefault}
\setstretch{1.15}

% Section formatting
\titleformat{\section}{\large\bfseries}{}{0em}{}[\vspace{-4pt}\rule{\linewidth}{0.4pt}\vspace{2pt}]
\titleformat{\subsection}{\normalsize\bfseries}{}{0em}{}
\titleformat{\subsubsection}{\normalsize\itshape}{}{0em}{}
\titlespacing{\section}{0pt}{18pt}{6pt}
\titlespacing{\subsection}{0pt}{10pt}{3pt}
\titlespacing{\subsubsection}{0pt}{6pt}{2pt}

% List settings
\setlist[itemize]{noitemsep, topsep=3pt, leftmargin=1.4em}
\setlist[enumerate]{noitemsep, topsep=3pt, leftmargin=1.4em}

% Table column types
\newcolumntype{L}[1]{>{\raggedright\arraybackslash}p{#1}}
\newcolumntype{R}[1]{>{\raggedleft\arraybackslash}p{#1}}

% Header / footer
\pagestyle{fancy}
\fancyhf{}
\renewcommand{\headrulewidth}{0.3pt}
\fancyhead[L]{\small\textbf{MedMemory}: Go-to-Market Note}
\fancyhead[R]{\small K V Naresh Karthigeyan}
\fancyfoot[C]{\small\thepage}

\begin{document}

% ─── TITLE BLOCK ─────────────────────────────────────────────────────────────
\begin{center}
  % {\large\textbf{MedMemory}}\\[4pt]
  {\large Go-to-Market Note}\\[6pt]
  {\large\textbf{MedMemory}}\\[4pt]
  % \rule{\linewidth}{0.6pt}\\[4pt]
  {\small K V Naresh Karthigeyan}
\end{center}

\vspace{3pt}
% ─── GTM ─────────────────────────────────────────────────────────────────────

\textit{Half a page, as requested. Sanity check, not a business plan.}

\subsection*{How the first 100 users find out}

The target early user is 28--35, lives in a metro, and manages health for a parent in a tier-2 city. They already feel the problem. They do not need convincing.

\begin{itemize}
  \item \textbf{WhatsApp communities.} Post in caregiver groups, NRI groups, chronic illness support circles. The hook is the WhatsApp upload flow --- it meets people where their report photos already live.
  \item \textbf{One LinkedIn post.} Not a product ad. A story: ``My mother has three doctors and I live in Bangalore. This is how I fixed it.'' Shares do the distribution.
  \item \textbf{Reddit and Quora.} r/india, r/IndianParents, r/ChronicIllness. Threads where people ask how to manage medical records for elderly parents. Answer genuinely, link as an aside.
  \item \textbf{Doctor waiting rooms.} Printed QR cards. Doctors are often the ones who wish patients would arrive better prepared.
\end{itemize}

\subsection*{What to measure to know they're actually using it}

Signups are noise. The metric that matters is \textbf{second upload rate}: what percentage of users who upload one report come back and upload a second within 30 days. That number tells you whether the product delivered enough value on the first use to earn the second.

\vspace{4pt}
\begin{tabularx}{\linewidth}{@{}L{7cm} X@{}}
\toprule
\textbf{Signal} & \textbf{What it tells you} \\
\midrule
Second upload rate (target: $>$40\% in 30 days) & Did the first experience create a reason to return? \\
Doctor brief generated per report uploaded & Are users reaching the output, or dropping at the record stage? \\
Verification rate on low-confidence entities & Are users trusting the AI enough to engage with corrections? \\
Emergency page views & High = people using it for its core promise, not just filing. \\
\midrule
\textit{Metric to explicitly ignore: DAU} & \textit{This is not a daily-use product. Monthly active is the right unit.} \\
\bottomrule
\end{tabularx}

\vspace{8pt}

One thing that would destroy early retention: a wrong medicine in the active list that a user does not catch. This is why human verification is in the critical path, not optional. The AI is wrong sometimes. The product has to be right.

\end{document}
